\documentclass[letterpaper,twocolumn]{article}  % Use twocolumn option
\usepackage[utf8]{inputenc}
\usepackage{geometry}              % Optional: for margin control
\usepackage{enumitem}
\geometry{margin=1in}     % Adjust margins if needed

\title{Calculating Energy Balances in Heat Exchanger Tests}
\author{Mike Stephenson, Solverware}
\date{March 2025}                  % Optional: customize date

\begin{document}
	
	\maketitle                         % Creates title section
	
	\begin{abstract}                   % Optional abstract
		Heat Exchanger performance testing involves setting up an experimental test rig, instrumenting it to measure temperatures, pressures and mass flow rates, acquiring the data, usually at several pre-defined state points and then evaluating the thermal effectiveness and pressure losses of the heat exchanger.
		
		It is important to calculate energy balances as one means of checking the integrity of the data. If there are leaks, inadequate insulation or other factors, the data may be suspect. An energy balance can help spot these factors.
	\end{abstract}
	
	\section{Introduction}
	When calculating an energy balance for a heat exchanger, you are checking to see how closely the heat lost by the hot fluid is equal to the heat gained by the cold fluid. There will be some system losses and other factors that will result in some deviation from the ideal. 
	
	\begin{equation}
		\dot{Q}_{cold} = \dot{Q}_{hot}
		\label{eq:darcy_equation}
	\end{equation}
	
	The well-known equation to calculate heat flux in a heat exchanger requires the mass flow rate, the specific heat and the temperature change. 
	
	\begin{equation}
		\dot{Q}=\dot{m}  C_p  \Delta T
	\end{equation}
	
	We calculate this on both the hot and cold side of the heat exchanger and compare.
	
	\begin{equation}
		\dot{Q}_{cold}=\dot{m}_{c}  C_{p,c} (T_{cold,out}-T_{cold,in})  
	\end{equation}
	
	\begin{equation}
		\dot{Q}_{hot}=\dot{m}_{h}  C_{p,h} (T_{hot,out}-T_{hot,in})  
	\end{equation}
	
	In an ideal heat exchanger (no losses), the heat gained by the cold fluid is the same as heat lost by the hot fluid. Mass flow rates and temperatures are measured by instrumentation. The specific heat, however, is usually a table lookup. This can get more complicated if the specific heat is a function of temperature which it often is. For these cases, we have some possible options:
	
	\begin{enumerate}[label=(\roman*)]
		\item Use the specific heat at the average temperature of the fluid stream.
		
		\item Look up the specific heat at both inlet and outlet temperatures and take the average
		
		\item Technically, the most accurate approach is to use the enthalpy difference
		
		\begin{equation}
			\dot{Q}_c=\dot{m}_c(h_{c,out}-h_{c,in})
		\end{equation}
	\end{enumerate}
	
	\section{The Enthalpy Difference Method}
	The enthalpy accounts for the specific heat's temperature dependence automatically.  For a constant $C_p$, the equation simplifies to what you often see in textbooks, i.e., $\dot{m}_cC_p\Delta T $. When specific heat has a temperature dependence, you would integrate this expression:
	
	\begin{equation}		
		h_{c,out}-h_{c,in} = \int_{T_{c,in}}^{T_{c,out}}C_p(T)dT
	\end{equation}
	
	Averaging $C_p$ at the inlet and outlet is essentially a shortcut to approximate this integral.
	
	In the majority of cases, this is usually good enough. But if your system has a fluid where the specific heat has a strong dependence on temperature, this is the way to go.
	
	You probably don't have an equation that provides the heat capacity as a function of temperature but if you do, it's a simple (or not so simple) matter of integration. If you have data points at a fixed number of equally-spaced temperatures over the integration range, you can numerically integrate to get the enthalpy change.
	
	Here's how you would do that:
	
	\begin{enumerate}
		\item Define the temperature range and points
			\begin{itemize}
				\item Pick a number of intervals, (n), to divide the range $T_{c,in}$ to $T_{c,out}$. This gives you $n+1$ points.
				\item Compute the temperature step size: 
				$$\Delta T = \frac{T_{c,out} - T_{c,in}}{n}$$
				\item Points: $T_i= T_{c,in} + i\cdot\Delta T$
			\end{itemize}
		\item Get $C_p$ at each point
			\begin{itemize}
				\item If you have routines that accomplish this, then simply write the code to get the $C_p$ as a function of $T$. Otherwise, look up the values from a reliable source.
			\end{itemize}	
		\item Choose an integration method
		
		For equally-spaced points, two common numerical integration methods are:
			\begin{enumerate}
				\item Trapezoidal Rule - Approximates the area under the curve by trapezoids. It’s straightforward and decently accurate for smooth $C_p(T)$ functions.
				\item Simpson's 1/3 Rule (if n is even) - Fits parabolas over pairs of intervals, usually more accurate but requires an even number of intervals (odd number of points).
			\end{enumerate}
		\item Computer the integral
		\item Calculate $\dot Q$ - 
		Once you have $\Delta H$, 
		compute $$\dot Q = m_c\cdot\Delta H $$
		
		
	\end{enumerate}
	
	
	\section{Main Content}
	
\end{document}